% TOP_PROBLEM: {"problemId": "problem.normality-of-pi", "canonicalTitle": "Normality of Pi", "releaseRank": 471, "primaryDomain": "number_theory_arithmetic_geometry", "primaryDomainLabel": "Number theory & arithmetic geometry", "displayStatus": "open", "releaseStatus": "open"}
% !TeX program = lualatex
% Definition and sources only; this file is not an LLM solution attempt.
\documentclass[11pt]{article}
\usepackage[margin=25mm]{geometry}
\usepackage{fontspec}
\setmainfont{Latin Modern Roman}
\usepackage{amsmath,amssymb,mathtools}
\usepackage{unicode-math}
\setmathfont{Latin Modern Math}
\usepackage{xurl,hyperref}
\hypersetup{colorlinks=true,urlcolor=blue}
\setcounter{secnumdepth}{0}
\setlength{\parindent}{0pt}
\setlength{\parskip}{5pt}
\setlength{\emergencystretch}{3em}
\begin{document}
\section*{\texorpdfstring{Normality of Pi}{Normality of Pi}}\label{471}
\subsection{Definitions and mathematical statement}
For an integer base $b\ge2$, write the fractional part of $\pi$ as
\[
 \pi-\lfloor\pi\rfloor=\sum_{j=1}^{\infty}a_jb^{-j},
 \qquad a_j\in\{0,\ldots,b-1\}.
\]
For a word $w=(w_1,\ldots,w_k)$ of length $k$, count its overlapping occurrences starting at positions $1$ through $N$. Normality in base $b$ means
\[
 \lim_{N\to\infty}\frac1N
 \#\{1\le j\le N:(a_j,\ldots,a_{j+k-1})=w\}=b^{-k}
\]
for every $k$ and every word $w$. The selected conjecture asks for this in every integer base, often called absolute normality. Equal frequencies for individual digits alone are insufficient, as are finite computations of very many digits.
\par\smallskip\textit{Formulation and concept references:} \href{https://www.claymath.org/library/annual_report/ar2006/06report_normalnumbers.pdf}{[S1]}.

\subsection{Short English statement}
In every integer base, do all finite digit patterns appear in pi with exactly the frequencies expected from uniform random digits?
\subsection{Sources}
\begin{itemize}
\item[S1]Davar Khoshnevisan, "Normal Numbers are Normal," Clay Mathematics Institute Annual Report 2006, pp. 15, 27-31.
\url{https://www.claymath.org/library/annual_report/ar2006/06report_normalnumbers.pdf}.
\textit{Formulation source.}
\end{itemize}
\end{document}
